\documentclass[a4paper,12pt,twoside,openright,titlepage]{book}

%Additional packages
\usepackage[utf8]{inputenc}
\usepackage[T1]{fontenc}
\usepackage[dutch,english]{babel}
\usepackage{imakeidx}
\usepackage{syntonly}
\usepackage[official]{eurosym}
%\usepackage[graphicx]
\usepackage{graphicx}
\graphicspath{ {./images/} }
\usepackage{float}
\usepackage{xurl}
\usepackage{hyperref}
\hypersetup{colorlinks=true, linkcolor=blue, citecolor=blue, filecolor=blue, urlcolor=blue, pdftitle=, pdfauthor=, pdfsubject=, pdfkeywords=}
\usepackage{tabularx}
\usepackage[table]{xcolor} % Table colors
\usepackage{scrextend}
\addtokomafont{labelinglabel}{\sffamily}
\usepackage{listings}
\usepackage{adjustbox}
\usepackage{color}

% Define colors
\definecolor{ashgrey}{rgb}{0.7, 0.75, 0.71}

% Listing style
\lstset{
  backgroundcolor=\color{ashgrey}, % choose the background color; you must add \usepackage{color} or \usepackage{xcolor}; should come as last argument
  basicstyle=\footnotesize,        % the size of the fonts that are used for the code
  breakatwhitespace=true,          % sets if automatic breaks should only happen at whitespace
  breaklines=true,                 % sets automatic line breaking
  extendedchars=true,              % lets you use non-ASCII characters; for 8-bits encodings only, does not work with UTF-8
  frame=single,	                   % adds a frame around the code
  rulecolor=\color{black},         % if not set, the frame-color may be changed on line-breaks within not-black text (e.g. comments (green here))
  keepspaces=true,                 % keeps spaces in text, useful for keeping indentation of code (possibly needs columns=flexible)
  columns=fullflexible,		   % make copy and paste possible
  showstringspaces=false,          % if true show spaces in strings adding particular underscores
  showspaces=false,                % if true show spaces everywhere adding particular underscores; it does not override 'showstringspaces'
}

% Uncomment for production
% \syntaxonly

% Style
\pagestyle{headings}

% Turn on indexing
\makeindex[intoc]

% Define document
\author{D. Leeuw}
\title{Hexadecimaal}
\date{\today\\v.0.8.0}

\begin{document}
\selectlanguage{dutch}

\maketitle

\copyright\ 2023 Dennis Leeuw\\

\begin{figure}
\includegraphics[width=0.3\textwidth]{CC-BY-SA-NC.png}
\end{figure}

\bigskip

Dit werk is uitgegeven onder de Creative Commons BY-NC-SA Licentie en laat anderen toe het werk te kopi\"eren, distribueren, vertonen, op te voeren, en om afgeleid materiaal te maken, zolang de auteurs en uitgever worden vermeld als maker van het werk, het werk niet commercieel gebruikt wordt en afgeleide werken onder identieke voorwaarden worden verspreid.


%%%%%%%%%%%%%%%%%%%
%%% Introductie %%%
%%%%%%%%%%%%%%%%%%%

\frontmatter
\chapter{Over dit Document}
\section{Leerdoelen}
Dit document geeft een basis uitleg over de boolean algebra in relatie tot de elektrotechniek. De verschillende Boolean functies worden beschreven aan de hand van elektronische schema's. De lezer leert dan ook wat boolean algebra is in de context van de elektrotechniek.


\section{Voorkennis}
\begin{itemize}
\item Om de uitleg te kunnen volgen die beschreven is bij de elektronische werking van de verschillende functies dient de lezer kennis te hebben van de werking van transistoren, weerstanden, spanning en stroom
\item Het is handig als de lezer enigzins vertrouwd is met het lezen van stroomschema's
\end{itemize}


%%%%%%%%%%%%%%%%%
%%% De inhoud %%%
%%%%%%%%%%%%%%%%%
\tableofcontents

\mainmatter
\chapter{Wat is hexadecimaal}
Het decimale systeem past niet echt lekker bij computers. Door het gebruik van enen en nullen in bits en bytes valt de decimale scheiding niet samen met de binaire scheiding van getallen. Een systeem dat hier beter bij past is het hexadecimale stelsel. De binaire reeks 1, 2, 4, 8, 16, 32, 64 past naadloos in een hexadecimale telling waarbij er geteld wordt van 0 tot f. Zoals de naam hexadecimaal al aangeeft gaat het om een 16-talig stelsel. We tellen in hexadecimaal van 0 tot 9 en daarna van A tot F voordat er een voorloop 1 voor komt, daarna weer van 0 tot 9 en A tot F voordat de voorloop 1 een 2 wordt. Zijn we bij het voorloopgetal aangekomen bij 9 dan tellen we door naar A en zo verder tot F.

Door van 0 tot f te tellen wordt de notatie van binaire getallen kleiner en beter leesbaar. Een lange reeks van enen en nullen wordt al snel onoverzichtelijk.

Door een stelsel te hebben waarbij er geteld wordt van 0 tot 16 is dit gelijk aan 4 bits in een binair stelsel.
Als we een nibble (4-bits) optellen (decimaal) komen we tot
\begin{math}
2^3 + 2^2 + 2^1 + 2^0 = 8 + 4 + 2 + 1 = 15
\end{math}
. We tellen daarmee van 0 tot 15 en dat is 16-talig. Wat we kunnen weergeven als 0-9 en A tot F.

4-bits heet een nibble. Toch werken we vaker met bytes, oftewel 8-bits. We hebben dus twee hexadecimale getallen nodig om een byte weer te geven, bijvoorbeel EA. Binnen een byte tellen we van 00 naar FF hexadecimaal wat decimaal overeen komt met 0 tot 255.



%%%%%%%%%%%%%%%%%%%%%
%%% Index and End %%%
%%%%%%%%%%%%%%%%%%%%%
\backmatter
\printindex
\end{document}

%%% Last line %%%
